\documentclass[../geometrie.tex]{subfiles}

\begin{document}
	\tableofcontents{}

	\section{Definition}{
		\subsection{Ensemble des \texorpdfstring{$\mathscr{C}^{\infty}$-fonctions}{C-inf fonctions}}{
			Soit $(M, \mathcal{A})$ une variété lisse.

			\mathblock{definition}{Ensemble $\mathscr{C}^{\infty}(M)$}{
				On note $\mathscr{C}^{\infty}(M)$ l'ensemble des $\mathscr{C}^{\infty}$-fonctions.
			}

			\mathblock{proposition}{Structure de $\mathscr{C}^{\infty}(M)$}{
				$\mathscr{C}^{\infty}$ définie une $\mathbb{R}$-algèbre.
			}
		}

		\subsection{Vecteur tangent, espace tangent}{
			Soit $(M, \mathcal{A})$ une variété lisse. \\
			Soit $p \in M$.

			\mathblock{definition}{Vecteur tangent}{
				Soit $v : \mathscr{C}^{\infty}(M) \longrightarrow \mathbb{R}$. \bigskip \\
				On dit que $v$ est un \textit{vecteur tangent} de $M$ en $p$ si
				\begin{mathlist}
					\item $v$ est linéaire.
					\item $v$ vérifie la relation de \textsc{Leibniz}.
						$$
							\forall f, g \in \mathscr{C}^{\infty}(M), v(fg) = f(p) v(g) + v(f) g(p)
						$$
				\end{mathlist}
			}

			\mathblock{definition}{Ensemble $\mathrm{T}_{p} M$}{
				On appelle \textit{ensemble des vecteurs tangents} de $M$ en $p$ et on note $\mathrm{T}_{p} M$ l'ensemble définie par
				$$
					\mathrm{T}_{p} M := \{ v \ \textit{vecteur tangent}  \}
				$$
			}

			\mathblock{proposition}{Structure de $\mathrm{T}_{p} M$}{
				$\mathrm{T}_{p} M$ est un $\mathbb{R}$-espace vectoriel.
			}

			\mathblock{lemma}{Vecteur tangent et fonction constante}{
				Soit $v \in \mathrm{T}_{p} M$. \\
				Soit $f \in \mathscr{C}^{\infty}(M)$ une fonction constante. \bigskip \\
				Alors
				$$
					v(f) = 0
				$$
			}

			\mathblock{lemma}{Vecteur tangent et fonction en $p$}{
				Soit $v \in \mathrm{T}_{p} M$. \\
				Soient $f, g \in \mathscr{C}^{\infty}(M)$.
				Supposons $f(p) = g(p) = 0$. \bigskip \\
				Alors
				$$
					v(fg) = 0
				$$
			}

			\mathblock{lemma}{Vecteur tangent en fonctions localement égales}{
				Soit $v \in \mathrm{T}_{p} M$. \\
				Soient $f, g \in \mathscr{C}^{\infty}(M)$. \\
				Supposons $\exists U \in \mathcal{O}_{p}$ tel que $f_{| U} = g_{| U}$. \bigskip \\
				Alors
				$$
					v(f) = v(g)
				$$
			}
		}

		\subsection{Espace tangent de \texorpdfstring{$\mathbb{R}^{n}$}{Rn}}{
			Soit $U \subset \mathbb{R}^{n}$ un ouvert. \\
			Soit $(e_{1}, \cdots, e_{n})$ la base canonique de $\mathbb{R}^{n}$. \\
			Soit $p \in U$. \\
			Soit $r \in \mathbb{R}^{n}$.

			\mathblock{definition}{Opérateur $\mathrm{D}_{p, r}$}{
				On appelle opérateur dérivée partielle sur $U$ en $p$ la fonctionnelle notée $\mathrm{D}_{p, r}$ et définie par
				$$
					\mathrm{D}_{p, r} :
					\begin{array}{rcl}
						\mathscr{C}^{\infty}(U) & \longrightarrow & \mathbb{R} \\
						f & \longmapsto & \displaystyle \frac{\partial f}{\partial r}(p)
					\end{array}
				$$
			}

			\mathblock{proposition}{$\mathrm{D}_{p, r}$ est un vecteur tangent}{
				$\mathrm{D}_{p, r} \in \mathrm{T}_{p} U$
			}

			\mathblock{theorem}{$\mathbb{R}^{n} \cong \mathrm{T}_{p} U$}{
				Soit l'application
				$$
					\mathrm{D}_{p} :
					\begin{array}{rcl}
						\mathbb{R}^{n} & \longrightarrow & \mathrm{T}_{p} U \\
						r & \longmapsto & \mathrm{D}_{p, r}
					\end{array}
				$$
				Alors $\mathrm{D}_{p}$ est un isomorphisme.
			}

			\mathblock{corollary}{Une base de $\mathrm{T}_{p} U$}{
				Soit $\mathcal{B} := (\mathrm{D}_{p, e_{1}}, \cdots, \mathrm{D}_{p, e_{n}})$. \bigskip \\
				Alors $\mathcal{B}$ est une base de $\mathrm{T}_{p} U$.
			}

			\mathblock{definition}{Opérateur dérivée partielle dans $\mathbb{R}^{n}$}{
				Soit $(M, \mathcal{A})$ une variété lisse de dimension $n$. \\
				Soit $(V, \varphi) \in \mathcal{A}$ telle que $\varphi(V) = U$. \\
				Soit $i \in [\![1, n]\!]$. \bigskip \\
				On appelle \textit{opérateur dérivée partielle} $U$ en $p$ selon la direction $i$ le vecteur tangent définie par
				$$
					\left. \frac{\partial}{\partial \varphi_{i}} \right|_{\varphi(p)} := \mathrm{D}_{p, e_{i}}
				$$
			}
		}
	}

	\section{Différentielle}{
		\noindent Soit $(M, \mathcal{A})$ une variété lisse. \\
		\noindent Soit $(N, \mathcal{B})$ une variété lisse. \\
		\noindent Soit $(Q, \mathcal{C})$ une variété lisse. \\
		\noindent Soit $F : M \longrightarrow N$ une $\mathscr{C}^{\infty}$-application. \\
		\noindent Soit $G : N \longrightarrow Q$ une $\mathscr{C}^{\infty}$-application. \\
		\noindent Soit $p \in M$.

		\mathblock{definition}{Différentielle de $F$ en $p$}{
			On appelle \textit{différentielle} de $F$ en $p$ l'application notée $\mathrm{d} F_{p}$ définie par
			$$
				\mathrm{d} F_{p} :
				\begin{array}{rcl}
					\mathrm{T}_{p} M & \longrightarrow & \mathrm{T}_{F(p)} N \\
					v & \longmapsto &
					\begin{array}{rcl}
						\mathscr{C}^{\infty}(N) & \longrightarrow & \mathbb{R} \\
						\varphi & \longmapsto & v(\varphi \circ F)
					\end{array}
				\end{array}
			$$
		}

		\mathblock{proposition}{$\mathrm{d} F_{p}$ est bien définie}{
			Soit $v \in \mathrm{T}_{p} M$. \bigskip \\
			Alors
			$$
				\mathrm{d} F_{p}(v) \in \mathrm{T}_{F(p)} N
			$$
		}

		\mathblock{proposition}{Composition de différentielles}{
			$$
				\mathrm{d}(G \circ F)_{p} = \mathrm{d} G_{F(p)} \circ \mathrm{d} F_{p}
			$$
		}

		\mathblock{lemma}{Différentielle de l'identité}{
			$$
				\mathrm{d} \mathrm{id}_{M} = \mathrm{id}_{\mathrm{T}_{p} M}
			$$
		}

		\mathblock{proposition}{$\mathrm{d} F_{p}$ est un isomorphisme}{
			$$
				\mathrm{d} F_{p} \in \mathrm{GL}(\mathrm{T}_{p} M, \mathrm{T}_{F(p)} N)
			$$
		}

		\mathblock{theorem}{Dimension de $\mathrm{T}_{p} M$}{
			$$
				\dim \mathrm{T}_{p} M = \dim M
			$$
		}
	}

	\section{Dérivée partielle dans une variété}{
		\noindent Soit $(M, \mathcal{A})$ une variété lisse de dimension $n$. \\
		\noindent Soit $(e_{1}, \cdots, e_{n})$ la base canonique de $\mathbb{R}^{n}$. \\
		\noindent Soit $p \in M$.

		\mathblock{definition}{Dérivée partielle dans $M$}{
			Soit $(U, \varphi) \in \mathcal{A}$ telle que $p \in U$. \\
			Soit $i \in [\![ 1, n]\!]$. \bigskip \\
			On appelle \textit{dérivée partielle selon $i$} en $p$ la dans $U$ fonctionnelle définie par
			$$
				\left. \frac{\partial}{\partial \varphi_{i}} \right|_{p} :
				\begin{array}{rcl}
					\mathscr{C}^{\infty}(U) & \longrightarrow & \mathbb{R} \\
					f & \longmapsto & \mathrm{D}_{\varphi(p), e_{i}}(f \circ \varphi^{-1})
				\end{array}
			$$
		}

		Dans le contexte de la $\mathscr{C}^{\infty}$-carte $(U, \varphi)$, les $n$ applications pour $i \in [\![1, n]\!]$
		$$
			\varphi_{i} : U \longrightarrow \mathbb{R}
		$$
		définissent les applications coordonnées. \\
		Si nous considèrons la configuration particulière de $(\mathbb{R}^{n}, \mathrm{id}_{\mathbb{R}^{n}})$, retrouvons en particulier les dérivées partielles de $\mathbb{R}^{n}$ en $p$ selon les directions canoniques
		$$
			\forall f \in \mathscr{C}^{\infty}(\mathbb{R}^{n}), \left. \frac{\partial}{\partial (\mathrm{id}_{\mathbb{R}^{n}})_{i}} \right|_{p}(f) = \frac{\partial f}{\partial x_{i}}(p)
		$$
		On remarque directement par définition de la différentielle que
		$$
			\left. \frac{\partial}{\partial \varphi_{i}} \right|_{p} = \mathrm{d}(\varphi^{-1})_{\varphi(d)}(\mathrm{D}_{\varphi(p), e_{i}})
		$$
		Il est également possible de définir une dérivée partielle sur l'ouvert $\varphi(U)$ dans $\mathbb{R}^{n}$.

		\mathblock{definition}{Dérivée partielle dans $\varphi(U)$}{
			Soit $(U, \varphi) \in \mathcal{A}$ telle que $p \in U$. \\
			Soit $i \in [\![ 1, n]\!]$. \bigskip \\
			On appelle \textit{dérivée partielle selon $i$} dans en $p$ dans $\varphi(U)$ la fonctionnelle définie par
			$$
				\left. \frac{\partial}{\partial \varphi_{i}} \right|_{\varphi(p)} :
				\begin{array}{rcl}
					\mathscr{C}^{\infty}(\varphi(U)) & \longrightarrow & \mathbb{R} \\
					f & \longmapsto & \mathrm{D}_{\varphi(p), e_{i}}(f)
				\end{array}
			$$
		}

		\mathblock{proposition}{Base de $\mathrm{T}_{p} U$}{
			On a
			$$
				\left( \left. \frac{\partial}{\partial \mathrm{\varphi}_{1}} \right|_{p}, \cdots, \left. \frac{\partial}{\partial \mathrm{id}_{n}} \right|_{p} \right)
			$$
			qui est une base de $\mathrm{T}_{p} U$.
		}

		En particulier n'importe quelle $\mathscr{C}^{\infty}$-carte permet de former une base dans son domaine.

		\mathblock{proof}{}{
			On a $\mathrm{d} (\varphi^{-1})_{\varphi(p)} : \mathrm{T}_{\varphi(p)} \varphi(U) \longrightarrow \mathrm{T}_{p} U$ qui est un isomorphisme. \\
			Or
			$$
				\left( \left. \frac{\partial}{\partial \varphi_{1}} \right|_{p}, \cdots \left. \frac{\partial}{\partial \varphi_{n}} \right|_{p} \right)
			$$
			est une base de $\mathrm{T}_{\varphi(p)} \varphi(U)$. \\
			De plus, pour $i \in [\![ 1, n  ]\!]$, et par la remarque précédente
			$$
				\mathrm{d} (\varphi^{-1})_{\varphi(p)} \! \left( \left. \frac{\partial}{\partial \varphi_{i}} \right|_{p} \right) = \left. \frac{\partial}{\partial \varphi_{i}} \right|_{p}
			$$
			D'où
			$$
				\left( \left. \frac{\partial}{\partial x_{1}} \right|_{p}, \cdots, \left. \frac{\partial}{\partial x_{n}} \right|_{p} \right)
			$$
			est une base de $\mathrm{T}_{p} U$.
		}

		\mathblock{lemma}{Dérivée partielle et application coordonnée}{
			On a
			$$
				\left. \frac{\partial}{\partial \varphi_{i}} \right|_{p}(\varphi_{k}) = \delta_{i, k}
			$$
		}

		\mathblock{proposition}{Décomposition}{
			Soient $(M, \mathcal{A}), (N, \mathcal{B})$ deux variété lisses de dimension $n$. \\
			Soit $p \in M$. \\
			Soit $F : M \longrightarrow N$ une $\mathscr{C}^{\infty}$-application. \\
			Soit $(U, \varphi) \in \mathcal{A}$ telle que $p \in U$. \\
			Soit $(V, \psi) \in \mathcal{B}$ telle que $F(U) \subset V$. \bigskip \\
			Alors
			$$
				\mathrm{d} F_{p} \left( \left. \frac{\partial}{\partial \varphi_{i}} \right|_{p} \right) = \sum_{j = 1}^{n} \left( \left. \frac{\partial}{\partial \varphi_{i}} \right|_{p} (\psi_{j} \circ F) \right) \times \left. \frac{\partial}{\partial \psi_{j}} \right|_{F(p)}
			$$
		}

		\mathblock{proof}{}{
			On remarque d'une part que
			$$
				\mathrm{d} F_{p} \! \left( \left. \frac{\partial}{\partial \varphi_{i}} \right|_{p} \right) \in \mathrm{T}_{F(p) N}
			$$
			Or
			$$
				\left( \left. \frac{\partial}{\partial \psi_{1}} \right|_{F(p)}, \cdots, \left. \frac{\partial}{\partial \psi_{n}} \right|_{F(p)} \right)
			$$
			est une base de $T_{F(p)} N$. \\
			Soit $\lambda_{1}, \cdots, \lambda_{n} \in \mathbb{R}$ des coefficients tels que
			$$
				\mathrm{d} F_{p} \! \left( \left. \frac{\partial}{\partial \varphi_{i}} \right|_{p} \right) = \sum_{j = 1}^{n} \lambda_{j} \times \left. \frac{\partial}{\partial \psi_{i}} \right|_{F(p)}
			$$
			Soit $k \in [\![1, n]\!]$. \\
			Alors par le lemme précédent,
			$$
				\sum_{j = 1}^{n} \lambda_{j} \times \left. \frac{\partial}{\partial \psi_{i}} \right|_{F(p)}(\psi_{k}) = \lambda_{k}
			$$
			Et on a par définition de la différentielle,
			$$
				\lambda_{k} = \mathrm{d} F_{p} \! \left( \left. \frac{\partial}{\partial \varphi_{i}} \right|_{p} \right)(\psi_{k}) = \left. \frac{\partial}{\partial \varphi_{i}} \right|_{p}(\psi_{k} \circ F)
			$$
			Finalement on retrouve bien la formule.
		}

		\mathblock{corollary}{Décomposition dans le cas euclidien}{
			
		}

		\mathblock{proposition}{Dérivée partielle et changement de carte}{
			On a
			$$
				\left. \frac{\partial}{\partial \varphi_{i}} \right|_{p} = \sum_{j = 1}^{n} \left. \frac{\partial}{\partial \varphi_{i}} \right|_{\varphi(p)}(\psi_{j} \circ \varphi^{-1}) \times \left. \frac{\partial}{\partial \psi_{j}} \right|_{p}
			$$
		}

		\mathblock{proof}{}{

		}
	}

	\section{Espace fibré tangent total}{
		Soit $(M, \mathcal{A})$ une variété lisse de dimension $n$.

		\subsection{Définition}{

			\mathblock{definition}{Espace tangent total}{
				On appelle \textit{espace tangent total de $M$} l'espace noté $\mathrm{T} M$ et définie par
				$$
					\begin{aligned}
						\mathrm{T} M & := \bigsqcup_{p \in M} {p} \times \mathrm{T}_{p} M \\
						& := \{ (p, v), p \in M, v \in \mathrm{T}_{p} M \}
					\end{aligned}
				$$
			}

			\mathblock{definition}{Projection naturelle}{
				On appelle \textit{projection naturelle de $\mathrm{T} M$ vers $M$} l'application notée $\pi$ définie par
				$$
					\pi :
					\begin{array}{rcl}
						\mathrm{T} M & \longrightarrow & M \\
						(p, v) & \longmapsto & p
					\end{array}
				$$
			}

			\mathblock{proposition}{Ensemble $\pi^{-1}(U)$}{
				$$
					\pi^{-1}(U) = \mathrm{T} U
				$$
			}

			\mathblock{proof}{}{
				\begin{itemize}
					\item[$\bullet$] Soit $(p, v) \in \mathrm{T} U$. \\
					Alors $p = \pi(p, v) \in U$. \\
					Donc $(p, v) \in \pi^{-1}(U)$. \\
					D'où $\mathrm{T} U \subset \pi^{-1}(U)$.

					\item[$\bullet$] Inversement, soit $(p, v) \in \pi^{-1}(U)$. \\
					Alors $p = \pi(p, v) \in U$. \\
					Donc $(p, v) \in \{ p \} \times \mathrm{T}_{p} U$.
				\end{itemize}
			}

			\mathblock{definition}{Extension $\tilde{\varphi}$}{
				Soit $(U, \varphi) \in \mathcal{A}$. \bigskip \\
				On appelle \textit{extension de $\varphi$} la fonction notée $\tilde{\varphi}$ définie par
				$$
					\tilde{\varphi} :
					\begin{array}{rcl}
						\pi^{-1}(U) & \longrightarrow & \varphi(U) \times \mathbb{R}^{n} \\
						\displaystyle \left(p, \sum_{i = 1}^{n} v_{i} \left. \frac{\partial}{\partial \varphi_{i}} \right|_{p} \right) & \longmapsto & \displaystyle \left(\varphi(p), \sum_{i = 1}^{n} v_{i} e_{i} \right)
					\end{array}
				$$
			}

			On vérifie que cette fonction est bien définie.

			\mathblock{proposition}{$\tilde{\varphi}$ est bijectif}{
				\begin{mathlist}
					\item $\tilde{\varphi}$ est bien définie.
					\item $\tilde{\varphi}$ est bijectif.
				\end{mathlist}
			}

			\mathblock{proof}{}{
				\begin{mathlist}
					\item Montrons que $\tilde{\varphi}$ est bien définie. \\
					Soient 
				\end{mathlist}
			}
		}

		\subsection{Topologie sur \texorpdfstring{$\mathrm{T} M$}{TM}}{
			\mathblock{proposition}{Topologie à partie d'une carte}{
				Soit $(U_{\alpha}, \varphi_{\alpha})_{\alpha \in A}$ une base dénombrable de $\mathscr{C}^{\infty}$-cartes sur M. \\
				Soit
				$$
					\begin{aligned}
						\mathcal{T}_{\varphi_{\alpha}} := & \{ \tilde{\varphi_{\alpha}}^{-1}(W), W \subset \varphi_{\alpha}(U_{\alpha}) \times \mathbb{R}^{n}, \\
						& W \ \text{ouvert pour} \ \varphi_{\alpha}(U_{\alpha}) \times \mathbb{R}^{n} \}
					\end{aligned}
				$$
				Alors
				\begin{mathlist}
					\item $\mathcal{T}_{\varphi_{\alpha}}$ définie une topologie sur $\mathrm{T} U_{\alpha}$.
					\item $\tilde{\varphi_{\alpha}} : \mathrm{T} U_{\alpha} \longrightarrow \varphi_{\alpha}(U_{\alpha}) \times \mathbb{R}^{n}$ est un homéomorphisme.
				\end{mathlist}
			}
		}
	}
\end{document}
